% Môn: Vật lý
% Lớp: 11
% Chương: Chương I. Dao động
% Bài: L11C1 Bài 2. Mô tả dao động điều hòa · Xác định chiều và tính chất chuyển động của vật dao động điều hòa
% ===== Câu 46 =====
% ID: L11C1B2-04-DS
% Mức: TH
\dangbt{Xác định chiều và tính chất chuyển động của vật dao động điều hòa}
\begin{ex}
% ID: L11C1B2-04-DS
% Mức: TH
Từ $x=7\cos(4t+\pi/3)$ cm, lập $v(t),a(t)$ và tính tại $t=0$.
\choiceTF

\end{ex}

% ===== Câu 50 =====
% ID: L11C1B2-05-DS
% Mức: VD
\begin{ex}
% ID: L11C1B2-05-DS
% Mức: VD
Vật có $A=8$ cm, $\omega=5$ rad/s, tại $x=A/2$ và đi theo chiều dương. Tính $v,a$ và xét nhanh dần hay chậm dần.
\choiceTF

\end{ex}

% ===== Câu 54 =====
% ID: L11C1B2-06-DS
% Mức: VD
\begin{ex}
% ID: L11C1B2-06-DS
% Mức: VD
Phân tích dấu của $x,v,a$ và sự thay đổi tốc độ khi vật đi từ biên âm về vị trí cân bằng.
\choiceTF

\end{ex}

% ===== Câu 61 =====
% ID: L11C1B2-07-DS
% Mức: TH
\begin{ex}
% ID: L11C1B2-07-DS
% Mức: TH
Vật dao động theo $x=6\cos(2t)$ cm.
\choiceTF
{\True $v=-12\sin(2t)$ cm/s.}
{\True $a=-24\cos(2t)$ cm/s².}
{\True $v_{max}=12$ cm/s.}
{Gia tốc bằng 0 ở biên.}
\end{ex}

% ===== Câu 65 =====
% ID: L11C1B2-08-DS
% Mức: VD
\begin{ex}
% ID: L11C1B2-08-DS
% Mức: VD
Vật có $A=7$ cm, $\omega=3$ rad/s, tại $x=A/2$ và đi theo chiều dương.
\choiceTF
{\True Gia tốc âm.}
{\True Độ lớn gia tốc bằng $31.5\,\text{cm}$ /s².}
{\True Vận tốc dương.}
{\True Vật chậm dần.}
\end{ex}

% ===== Câu 69 =====
% ID: L11C1B2-09-DS
% Mức: VD
\begin{ex}
% ID: L11C1B2-09-DS
% Mức: VD
Vật đi từ biên âm về vị trí cân bằng.
\choiceTF
{\True Li độ âm.}
{\True Vận tốc dương.}
{\True Gia tốc dương.}
{Tốc độ giảm dần.}
\end{ex}

% ===== Câu 71 =====
% ID: L11C1B2-10-TLN
% Mức: TH
\begin{ex}
% ID: L11C1B2-10-TLN
% Mức: TH
Với $x=9\cos(5t)$ cm, tốc độ cực đại theo cm/s bằng bao nhiêu?
\shortans{45}
\end{ex}

% ===== Câu 75 =====
% ID: L11C1B2-11-TLN
% Mức: VD
\begin{ex}
% ID: L11C1B2-11-TLN
% Mức: VD
Tại $x=2.5$ cm, $\omega=2$ rad/s. Độ lớn gia tốc theo cm/s² bằng bao nhiêu?
\shortans{10}
\end{ex}

% ===== Câu 79 =====
% ID: L11C1B2-12-TLN
% Mức: VD
\begin{ex}
% ID: L11C1B2-12-TLN
% Mức: VD
Vật có $v_{max}=18$ cm/s và $A=6$ cm. Tần số góc theo rad/s bằng bao nhiêu?
\shortans{3}
\end{ex}

% ===== Câu 144 =====
% ID: L11C1B2-13-TN
% Mức: TH
\begin{ex}
% ID: L11C1B2-13-TN
% Mức: TH
Trong dao động điều hoà, gia tốc biến đổi
\choice
{cùng pha với vận tốc.}
{\True sớm pha 900 so với vận tốc.}
{ngược pha với vận tốc.}
{trễ pha 900 so với vận tốc.}
\loigiai{
Trong dao động điều hoà, gia tốc biến đổi sớm pha 900 so với vận tốc.
}
\end{ex}

% ===== Câu 145 =====
% ID: L11C1B2-14-TN
% Mức: TH
\begin{ex}
% ID: L11C1B2-14-TN
% Mức: TH
Trong dao động điều hoà, vận tốc biến đổi
\choice
{ngược pha với gia tốc.}
{cùng pha với Li độ.}
{ngược pha với gia tốc.}
{\True sớm pha 900 so với Li độ.}
\loigiai{
Trong dao động điều hoà,vận tốc biến đổi sớm pha 900 so với Li độ.
}
\end{ex}

% ===== Câu 146 =====
% ID: L11C1B2-15-TN
% Mức: TH
\begin{ex}
% ID: L11C1B2-15-TN
% Mức: TH
Khi một vật dao động điều hòa, chuyển động của vật từ vị trí biên về vị trí cân bằng là chuyển động
\choice
{nhanh dần đều.}
{chậm dần đều.}
{\True nhanh dần.}
{chậm dần.}
\loigiai{
Khi một vật dao động điều hòa, chuyển động của vật từ vị trí biên về vị trí cân bằng là chuyển động nhanh dần.
}
\end{ex}

% ===== Câu 147 =====
% ID: L11C1B2-16-TN
% Mức: TH
\begin{ex}
% ID: L11C1B2-16-TN
% Mức: TH
Khi một vật dao động điều hòa, chuyển động của vật từ vị trí cân bằng ra vị trí biên dương là chuyển động
\choice
{nhanh dần đều.}
{chậm dần đều.}
{nhanh dần.}
{\True chậm dần.}
\loigiai{
Khi một vật dao động điều hòa, chuyển động của vật từ vị trí cân bằng ra vị trí biên dương là chuyển động chậm dần.
}
\end{ex}

% ===== Câu 148 =====
% ID: L11C1B2-17-TN
% Mức: TH
\begin{ex}
% ID: L11C1B2-17-TN
% Mức: TH
Khi một vật dao động điều hòa, chuyển động của vật từ vị trí cân bằng ra vị trí biên âm là chuyển động
\choice
{nhanh dần đều.}
{chậm dần đều.}
{nhanh dần.}
{\True chậm dần.}
\loigiai{
Khi một vật dao động điều hòa, chuyển động của vật từ vị trí cân bằng ra vị trí biên âm là chuyển động chậm dần.
}
\end{ex}
